\subsection{相交的超平面}

回忆（《代数》，第 II 章，§9，第 3 号），E 的两个仿射子空间 P 和 P' 称为平行的，如果存在 T 中的向量 t，使 \(P' = t + P\)。显然，关系“P 和 P' 平行”是 E 的仿射子空间集上的等价关系。

引理 1. 任意两个不平行的超平面都有非空交。

设 H 和 \(H'\) 是两个不平行的超平面，\(a \in H\) 且 \(a' \in H'\)；存在 T 这个向量空间的两个超平面 M 和 \(M'\)，使 \(H = M + a\) 且 \(H' = M' + a'\)；由于 H 和 \(H'\) 不平行，有 \(M \neq M'\)，因而 \(T = M + M'\)；于是存在 \(u \in M\) 和 \(u' \in M'\)，使 \(a' - a = u - u'\)，而点 \(u + a = u' + a'\) 属于 \(H \cap H'\)。

引理 2. 设 H 和 H' 是 E 的两个不同超平面，f 和 \(f'\) 是 E 上的两个仿射函数，使得 H（分别地，H'）由 E 中满足 \(f(a) = 0\)（分别为 \(f'(a) = 0\)）的点组成。最后，设 L 是 E 的一个超平面。假设下列假设之一成立：

\begin{enumerate}
\def\labelenumi{\alph{enumi})}
\item
  超平面 H、H' 和 L 平行。
\item
  超平面 H 和 \(\mathrm{H}'\) 不平行，且 \(\mathrm{H} \cap \mathrm{H}' \subset \mathrm{L}\)。
\end{enumerate}

那么存在不全为零的实数 \(\lambda, \lambda'\)，使得 \(L\) 由 \(a \in E\) 中使仿射函数 \(g = \lambda.f + \lambda'.f'\) 为零的那些点组成。

当 \(\mathbf{L} = \mathbf{H}\) 时引理显然成立，因此假设 \(a\) 中存在一点 \(\mathbf{L}\) 满足 \(a \notin \mathbf{H}\)。置 \(\lambda = f'(a)\)，\(\lambda' = -f(a)\)，并令

\[
g = \lambda . f + \lambda^ {\prime}. f ^ {\prime};
\]

于是由于 \(\lambda' \neq 0\)，有 \(a \notin \mathrm{H}\)；此外，由于 \(\mathrm{H} \neq \mathrm{H}'\)，存在 \(b \in \mathrm{H}\) 使 \(b \notin \mathrm{H}'\)，从而 \(f(b) = 0\)、\(f'(b) \neq 0\)，因而 \(g(b) = -f(a).f'(b)\) 非零。于是仿射函数 \(\mathbf{L}_1\) 为零的点集 \(g \neq 0\) 是 \(\mathbf{E}\) 的一个超平面；我们有 \(g(a) = 0\)，所以 \(a \in \mathbf{L}_1\)。

\begin{enumerate}
\def\labelenumi{\alph{enumi})}
\item
  假设 H 和 \(H'\) 平行：g 和 f 在 \(L_{1} \cap H\) 的每一点都为零，因此 \(f'\) 也如此，因为 \(\lambda' \neq 0\)；因此 \(L_{1} \cap H\) 的每一点都属于 \(H'\)；但由于 H 和 \(H'\) 平行且不同，它们不相交，所以 \(L_{1} \cap H = \varnothing\)，而引理 1 表明 \(L_{1}\) 与 H 平行。由于 \(a \in L\) 且 \(a \in L_{1}\)，故有 \(L = L_{1}\)。
\item
  假设 H 和 \(\mathbf{H}'\) 不平行：由引理 1，存在一点 \(c\) 属于 \(\mathbf{H} \cap \mathbf{H}'\)；以 \(c\) 为原点给 E 赋予向量空间结构。于是 \(H \cap H'\) 是 E 的余维为 2 的向量子空间，并且由于 \(a \notin H\)，由 \(H \cap H'\) 和 a 生成的 E 的向量子空间 M 是一个超平面；由于 \(H \cap H' \subset L \cap L_1\) 且 \(a \in L \cap L_1\)，有 \(M \subset L \cap L_1\)，从而 \(M = L = L_1\)。证毕。
\end{enumerate}

命题 10. 设 C 是一个室，H 和 H' 是 C 的两面墙，L 是与 \(D_{H}(C) \cap D_{H'}(C)\) 相交的超平面。假设 H 与 H' 不同，且下列条件之一成立：

\begin{enumerate}
\def\labelenumi{\alph{enumi})}
\item
  超平面 H、H' 和 L 平行。
\item
  超平面 H 和 H' 不平行，且 H ∩ H' ⊂ L。
\end{enumerate}

则 L 与 C 相交。

设 \(b\)（分别地，\(b'\)）是支撑为 H（分别为 H'）的 C 的面的一个点；立即可见，线段 \([bb']\) 中不同于 \(b\) 和 \(b'\) 的每一点都属于 C。

引入仿射函数 \(f\)，它在 \(H\) 的每一点都为零，并且对于 \(f(x) > 0\) 中的 \(x\) 满足 \(D_H(C)\)；同样，引入关于 \(f'\) 具有类似性质的仿射函数 \(H'\)。应用引理 2，可找到数 \(\lambda\) 和 \(\lambda'\) 以及具有引理中所述性质的仿射函数 \(g\)。有 \((\lambda, \lambda') \neq (0, 0)\)，并且对于每一点 \(x\)，若其属于 \(L \cap D_H(C) \cap D_{H'}(C)\)，有 \(f(x) > 0\)、\(f'(x) > 0\) 以及 \(\lambda.f(x) + \lambda'.f'(x) = 0\)，所以 \(\lambda\lambda' < 0\)。另一方面，有 \(g(b) = \lambda'.f'(b)\) 和 \(g(b') = \lambda.f(b')\)，由于 \(f(b') > 0\)、\(f'(b) > 0\)，有 \(g(b)g(b') < 0\)。于是点 \(b\) 和 \(b'\) 严格位于超平面 \(L\) 的相反两侧，因此存在一点 \(c\) 属于 \(L\)，它属于 \([bb']\)，且不同于 \(b\) 和 \(b'\)，从而属于 \(C\)。

